\chapter{第8章\texorpdfstring{\hspace{0.5em}}{ }电磁场与麦克斯韦电磁场方程组}

\section{电磁感应的基本定律}
    \subsection{电磁感应现象}
    当穿过闭合导体回路的磁通量发生变化时，回路中就会产生电流，这种现象称为电磁感应，产生的电流称为感应电流；回路中相应的电动势称为感应电动势。使磁通量发生变化的方式有多种：磁铁与线圈相对运动、线圈在磁场中运动或转动、线圈面积变化、磁场本身（电流）变化等。

    电磁感应现象是法拉第于 1831 年发现的，它揭示了电与磁之间的相互转化，是电磁理论的实验基础，也是发电机、变压器、电动机等一切电磁装置的工作原理。

    \subsection{楞次定律}
    \textbf{楞次定律}：闭合回路中感应电流的方向，总是使感应电流所产生的磁通量阻碍原磁通量的变化。

    楞次定律的实质是能量守恒定律在电磁感应中的体现：感应电流的磁场"阻碍"磁通量的变化，意味着必须由外界做功才能维持这种变化，从而把其他形式的能量转化为电能。判断感应电流方向的步骤为：（1）确定原磁场的方向；（2）判断穿过回路的磁通量是增加还是减少；（3）由楞次定律确定感应电流的磁场方向（增反减同）；（4）用右手螺旋法则由感应磁场方向确定感应电流方向。

    \subsection{法拉第电磁感应定律}
    精确的实验表明，回路中感应电动势的大小与穿过回路的磁通量对时间的变化率成正比。在国际单位制中，比例系数为 1：
    \keyformula{\varepsilon=-\frac{\mathrm{d}\Phi_m}{\mathrm{d}t}}
    式中负号反映了感应电动势的方向总是阻碍磁通量的变化，是楞次定律的数学表示。对 $N$ 匝紧密绕制的线圈，若每匝的磁通量均为 $\Phi_m$，则总电动势为
    \[
    \varepsilon=-N\frac{\mathrm{d}\Phi_m}{\mathrm{d}t}=-\frac{\mathrm{d}\Psi}{\mathrm{d}t},
    \]
    其中 $\Psi=N\Phi_m$ 称为磁通量链（磁链）。

    应当注意，$\Phi_m$ 是穿过以回路为边界的任意曲面的磁通量；$\varepsilon$ 的方向由楞次定律判定，式中负号只表示"阻碍"，具体方向要结合回路正方向的选取判断。若回路电阻为 $R$，则感应电流 $I=\varepsilon/R$，在 $\Delta t$ 时间内通过回路任一截面的电荷量为
    \[
    q=\int I\,\mathrm{d}t=\frac{1}{R}\left|\Delta\Phi_m\right|,
    \]
    电荷量只与磁通量的变化量有关，与变化的快慢无关，这是测量磁通量（冲击电流计）的原理。

    \corepoint{（1）法拉第定律中的 $\mathrm{d}\Phi_m/\mathrm{d}t$ 是磁通量对时间的变化率，磁通量大而恒定不产生感应电动势；（2）式中的负号体现楞次定律，不能漏掉，也不能据此直接断定电动势的绝对方向；（3）感应电动势的产生与回路是否闭合无关（闭合时有电流，不闭合时只在两端产生电势差）。}

\section{动生电动势}
    \subsection{动生电动势的产生原因}
    导体在磁场中运动（或导体回路面积变化）而产生的感应电动势称为动生电动势。以直导线在均匀磁场中平动为例：导线中的自由电子随导线以速度 $\boldsymbol{v}$ 运动，受到洛伦兹力
    \[
    \boldsymbol{F}_m=-e\boldsymbol{v}\times\boldsymbol{B},
    \]
    电子沿导线向一端聚集，另一端带正电，直到导线两端电荷产生的电场力与洛伦兹力平衡为止。可见，动生电动势的非静电力就是洛伦兹力，相应的非静电场强为
    \[
    \boldsymbol{E}_k=\frac{\boldsymbol{F}_m}{q}=\boldsymbol{v}\times\boldsymbol{B}.
    \]
    于是动生电动势为
    \keyformula{\varepsilon=\int_a^b\boldsymbol{E}_k\cdot\mathrm{d}\boldsymbol{l}=\int_a^b(\boldsymbol{v}\times\boldsymbol{B})\cdot\mathrm{d}\boldsymbol{l}}
    这个积分沿导体从 $a$ 到 $b$ 进行。

    \subsection{动生电动势的计算}
    对长为 $l$ 的直导线，在均匀磁场中以速度 $\boldsymbol{v}$ 平动，且 $\boldsymbol{v}$、$\boldsymbol{B}$、$\boldsymbol{l}$ 三者两两垂直时，
    \[
    \varepsilon=Blv.
    \]
    一般情形按 $\varepsilon=\int(\boldsymbol{v}\times\boldsymbol{B})\cdot\mathrm{d}\boldsymbol{l}$ 计算。动生电动势的方向可由 $\boldsymbol{v}\times\boldsymbol{B}$ 的方向判定（正电荷受力的方向即电动势方向，也就是从低电势指向高电势的方向）。

    计算动生电动势有两种常用方法：
    \begin{enumerate}
        \item 用 $\varepsilon=\int(\boldsymbol{v}\times\boldsymbol{B})\cdot\mathrm{d}\boldsymbol{l}$ 直接积分（适用于一段导线）；
        \item 用 $\varepsilon=-\dfrac{\mathrm{d}\Phi_m}{\mathrm{d}t}$ 求（适用于闭合回路，先算磁通量再求导）。
    \end{enumerate}
    由第（2）法可见，动生电动势也可看成是回路面积变化引起磁通量变化的结果，两种方法结果一致。

    \corepoint{动生电动势的非静电力是洛伦兹力，因此电动势的方向由 $\boldsymbol{v}\times\boldsymbol{B}$ 决定。注意：$\varepsilon=\int(\boldsymbol{v}\times\boldsymbol{B})\cdot\mathrm{d}\boldsymbol{l}$ 中若 $\boldsymbol{v}$ 与 $\boldsymbol{B}$ 平行，则 $\varepsilon=0$，即导线沿磁场方向运动时不产生动生电动势。}

\section{感生电动势和感生电场}
    \subsection{感生电动势}
    当导体不动而磁场随时间变化时，回路中同样产生感应电动势，这种电动势称为感生电动势。由于导体并未运动，洛伦兹力不能解释其来源。麦克斯韦指出，变化的磁场会在空间激发一种电场，称为感生电场（涡旋电场），用 $\boldsymbol{E}_i$ 表示。感生电场对导体中自由电荷的作用力就是产生感生电动势的非静电力：
    \[
    \varepsilon=\oint_L\boldsymbol{E}_i\cdot\mathrm{d}\boldsymbol{l}.
    \]

    \subsection{感生电场}
    由法拉第电磁感应定律与电动势的定义得
    \keyformula{\oint_L\boldsymbol{E}_i\cdot\mathrm{d}\boldsymbol{l}=-\frac{\mathrm{d}\Phi_m}{\mathrm{d}t}=-\frac{\mathrm{d}}{\mathrm{d}t}\int_S\boldsymbol{B}\cdot\mathrm{d}\boldsymbol{S}}
    这就是感生电场的环路定理。式中 $S$ 是以闭合路径 $L$ 为边界的任意曲面。

    感生电场与静电场有本质区别：
    \begin{table}[htbp]
    \centering
    \caption{静电场与感生电场的比较}
    \begin{tabular}{lll}
    \toprule
    项目 & 静电场 & 感生电场 \\
    \midrule
    场源 & 静止电荷 & 变化的磁场 \\
    电场线 & 起于正电荷、止于负电荷，不闭合 & 闭合曲线，无头无尾 \\
    环流 & $\oint\boldsymbol{E}\cdot\mathrm{d}\boldsymbol{l}=0$（保守场） & $\oint\boldsymbol{E}_i\cdot\mathrm{d}\boldsymbol{l}=-\mathrm{d}\Phi_m/\mathrm{d}t$（非保守场） \\
    电势 & 可引入电势 & 不能引入电势（非保守） \\
    \bottomrule
    \end{tabular}
    \end{table}

    感生电场是电磁相互作用的又一种表现，它不需要导体存在，即使在真空中也能存在；它的存在已被实验（如电子感应加速器）所证实。

    \subsection{感生电动势的计算}
    由 $\varepsilon=-\mathrm{d}\Phi_m/\mathrm{d}t$ 或 $\varepsilon=\oint\boldsymbol{E}_i\cdot\mathrm{d}\boldsymbol{l}$ 计算。对具有圆柱对称性的变化磁场（如半径为 $R$ 的长直螺线管内磁场均匀变化），感生电场线是与轴同心的圆，由环路定理可得其大小
    \[
    E_i=\frac{r}{2}\left|\frac{\mathrm{d}B}{\mathrm{d}t}\right|\quad(r<R),\qquad
    E_i=\frac{R^{2}}{2r}\left|\frac{\mathrm{d}B}{\mathrm{d}t}\right|\quad(r>R).
    \]
    可见管内 $E_i$ 与 $r$ 成正比，管外与 $r$ 成反比。

    \subsection{涡电流}
    当大块金属处在变化的磁场中时，金属内部也会产生感生电流，这种电流呈涡旋状，称为涡电流（涡流）。涡电流有两种效应：
    \begin{enumerate}
        \item 热效应：涡流在金属中产生焦耳热。可利用它进行感应加热（高频感应炉、真空冶炼）和金属探测；但在变压器、电机的铁芯中涡流发热是有害的，常采用互相绝缘的硅钢片叠压铁芯来减小涡流；
        \item 机械效应（电磁阻尼）：涡流在磁场中受安培力，其效果总是阻碍导体与磁场的相对运动。电磁阻尼可用于电表、天平的阻尼，电磁驱动（如感应式电动机）则利用它的反作用；
        \item 趋肤效应：高频电流几乎全部集中在导体表面附近流动，使导体有效截面减小、电阻增大，高频电路中常用多股绝缘细导线（辫线）或镀银导体来减小趋肤效应的影响。
    \end{enumerate}

    \corepoint{感生电场由变化的磁场激发，与是否存在导体无关；它不满足保守场条件，故不能引入电势。计算感生电动势时要注意磁场的对称性和"磁场变化率"的正负，方向仍由楞次定律判定。}

\section{自感和互感}
    \subsection{自感}
    当回路中的电流发生变化时，它所产生的磁场穿过回路自身的磁通量也随之变化，从而在回路中产生感应电动势，这种现象称为自感。由毕奥-萨伐尔定律，穿过回路的磁通量与电流成正比：
    \[
    \Psi=LI,
    \]
    比例系数 $L$ 称为回路的自感（自感系数），单位为亨利（$\mathrm{H}$），$1\ \mathrm{H}=1\ \mathrm{Wb/A}=1\ \mathrm{V\cdot s/A}$。自感只由回路的形状、大小、匝数以及周围磁介质的性质决定，与电流无关。

    由法拉第定律得自感电动势
    \keyformula{\varepsilon_L=-L\frac{\mathrm{d}I}{\mathrm{d}t}}
    负号表示自感电动势总是阻碍回路中电流的变化（"通电流时阻碍增大，断电流时阻碍减小"）。自感现象是电路中产生"电火花"、电流不能突变的原因。

    典型自感：长为 $l$、截面积为 $S$、单位长度匝数为 $n$ 的长直螺线管
    \[
    L=\mu_0\mu_rn^{2}V=\mu_0\mu_rn^{2}lS;
    \]
    同轴电缆（长 $l$，内外半径 $R_1$、$R_2$）
    \[
    L=\frac{\mu_0}{2\pi}l\ln\frac{R_2}{R_1}.
    \]

    \subsection{互感}
    两个相邻线圈，当其中一个线圈中的电流变化时，会在另一个线圈中产生感应电动势，这种现象称为互感。若线圈 1 中电流 $I_1$ 产生的磁场穿过线圈 2 的磁链为 $\Psi_{21}$，则
    \[
    \Psi_{21}=MI_1,
    \]
    $M$ 称互感系数，单位也是亨利。线圈 2 中的互感电动势为
    \[
    \varepsilon_{21}=-M\frac{\mathrm{d}I_1}{\mathrm{d}t},
    \]
    同理 $\varepsilon_{12}=-M\dfrac{\mathrm{d}I_2}{\mathrm{d}t}$。可以证明互感系数 $M_{12}=M_{21}=M$，即两个线圈的互感系数是相同的，它只由两线圈的形状、大小、相对位置和磁介质决定。

    两个共轴长直螺线管（长 $l$、截面积 $S$，单位长度匝数分别为 $n_1$、$n_2$）的互感为
    \[
    M=\mu_0\mu_rn_1n_2V=\mu_0\mu_rn_1n_2lS.
    \]
    一般情况下 $M=k\sqrt{L_1L_2}$，其中 $k$（$0\le k\le1$）称耦合系数，$k=1$ 时称全耦合（无漏磁）。

    互感有广泛的应用（变压器、互感器、无线电耦合电路），但在电子线路中相邻线圈的互感会引起相互干扰，常需采取屏蔽、垂直放置等措施减小互感。

    \corepoint{（1）自感 $L$ 与电流无关，只由回路本身决定；（2）自感电动势的方向由楞次定律判断，式中负号表示阻碍电流变化；（3）互感系数 $M$ 与两线圈的相对位置有关，$M\le\sqrt{L_1L_2}$。}

\section{RL 电路 磁场的能量}
    \subsection{RL 电路的暂态过程}
    含自感的电路（$R$、$L$ 串联）接通电源时，由于自感电动势的阻碍作用，电流不能立即达到稳定值。由基尔霍夫定律
    \[
    \varepsilon-L\frac{\mathrm{d}I}{\mathrm{d}t}=IR,
    \]
    在 $I(0)=0$ 的初始条件下解得
    \[
    I=\frac{\varepsilon}{R}\left(1-\mathrm{e}^{-\frac{R}{L}t}\right)=\frac{\varepsilon}{R}\left(1-\mathrm{e}^{-t/\tau}\right),
    \]
    其中 $\tau=L/R$ 称时间常数，它决定电流增长的快慢。经过 $t=\tau$ 时电流达到稳定值的 $63\%$，经过 $3\tau\sim5\tau$ 可近似认为达到稳定。

    断开电源时，电流按指数规律衰减：
    \[
    I=\frac{\varepsilon}{R}\mathrm{e}^{-\frac{R}{L}t}.
    \]
    由于电流突变会在 $L$ 上产生很大的自感电动势，实际电路中常在开关处产生火花，需加灭弧装置或续流二极管。

    \subsection{自感磁能}
    在电流由零增大到 $I$ 的过程中，电源克服自感电动势做功，把能量储存在线圈的磁场中。自感磁能为
    \keyformula{W_m=\frac{1}{2}LI^{2}}
    这与电容器的储能 $W_e=\dfrac{Q^{2}}{2C}$ 在形式上完全对应，说明磁场与电场一样具有能量。

    \subsection{磁场的能量 磁能密度}
    以长直螺线管为例，$L=\mu n^{2}V$，$B=\mu nI$，代入 $W_m=\frac{1}{2}LI^{2}$ 得
    \[
    W_m=\frac{1}{2}\mu n^{2}VI^{2}=\frac{B^{2}}{2\mu}V.
    \]
    于是单位体积内的磁场能量，即磁能密度为
    \keyformula{w_m=\frac{W_m}{V}=\frac{B^{2}}{2\mu_0\mu_r}=\frac{1}{2}BH=\frac{1}{2}\boldsymbol{B}\cdot\boldsymbol{H}}
    此式虽由螺线管导出，但对任意磁场都成立。任意磁场中的总能量为
    \[
    W_m=\int_Vw_m\,\mathrm{d}V=\int_V\frac{B^{2}}{2\mu}\,\mathrm{d}V.
    \]

    \corepoint{（1）磁能储存在磁场中，能量密度 $w_m=B^{2}/(2\mu)$，与电场能量密度 $w_e=\frac{1}{2}\varepsilon E^{2}$ 形式对称；（2）$RL$ 电路的时间常数 $\tau=L/R$，$L$ 越大、$R$ 越小，暂态过程越长；（3）断电时若强行断开电路，$\mathrm{d}I/\mathrm{d}t$ 很大，自感电动势可能远高于电源电压，这是"断电火花"的原因。}

\section{位移电流和全电流定律}
    \subsection{位移电流}
    恒定磁场的安培环路定理 $\oint\boldsymbol{H}\cdot\mathrm{d}\boldsymbol{l}=\sum I_0$ 在非恒定情况下遇到困难。考虑电容器充电电路：取一个以导线为边界的闭合路径 $L$，若以穿过导线的平面为积分面，则包围电流 $I$；若以穿过电容器两极板间的曲面为积分面，则没有传导电流穿过，$\sum I_0=0$。同一路径得到两个不同的结果，说明安培环路定理需要推广。

    麦克斯韦注意到，电容器充电时两极板间的电场在变化，电位移矢量 $\boldsymbol{D}$ 的通量也在变化。他提出：变化的电场等效于一种电流，称为位移电流，其定义为
    \keyformula{I_d=\frac{\mathrm{d}\Psi_D}{\mathrm{d}t}=\frac{\mathrm{d}}{\mathrm{d}t}\int_S\boldsymbol{D}\cdot\mathrm{d}\boldsymbol{S}=\int_S\frac{\partial\boldsymbol{D}}{\partial t}\cdot\mathrm{d}\boldsymbol{S}}
    相应的位移电流密度为
    \[
    \boldsymbol{j}_d=\frac{\partial\boldsymbol{D}}{\partial t}.
    \]
    在上述电容器电路中，导线中的传导电流 $I$ 与极板间的位移电流 $I_d$ 大小相等、方向相同（就电流的连续性而言），从而保证了电流的连续性。

    位移电流与传导电流有本质区别：位移电流不是电荷的定向移动，不产生焦耳热；它只存在于电场变化的空间（可以是真空、介质），其周围同样能激发磁场。位移电流的引入揭示了"变化的电场激发磁场"这一基本规律，与"变化的磁场激发电场"相对应。

    \subsection{全电流定律}
    传导电流与位移电流之和称为全电流。全电流是连续的。把安培环路定理推广到非恒定情形，得全电流定律
    \keyformula{\oint_L\boldsymbol{H}\cdot\mathrm{d}\boldsymbol{l}=I_0+I_d=\sum I_0+\int_S\frac{\partial\boldsymbol{D}}{\partial t}\cdot\mathrm{d}\boldsymbol{S}}
    即磁场强度沿任意闭合路径的环流，等于该路径所包围的传导电流与位移电流的代数和。

    \corepoint{（1）位移电流的本质是变化电场，不是电荷运动，不产生焦耳热；（2）$I_d=\mathrm{d}\Psi_D/\mathrm{d}t$ 与 $I=\mathrm{d}q/\mathrm{d}t$ 形式相似，但物理内容完全不同；（3）全电流定律中的 $\sum I_0$ 是穿过以 $L$ 为边界的曲面的传导电流，需注意曲面选取的一致性。}

\section{麦克斯韦电磁场方程组}
    \subsection{麦克斯韦方程组的积分形式}
    麦克斯韦在前人（库仑、高斯、安培、法拉第）成果的基础上，引入感生电场和位移电流两个基本概念，把电磁场的全部基本规律概括为四个方程，称为麦克斯韦电磁场方程组。其积分形式为
    \begin{align*}
    &\oint_S\boldsymbol{D}\cdot\mathrm{d}\boldsymbol{S}=\sum_{S\text{内}}q_{0i}
    &&\text{（电场的高斯定理：电场是有源场）}\\
    &\oint_L\boldsymbol{E}\cdot\mathrm{d}\boldsymbol{l}=-\frac{\mathrm{d}\Phi_m}{\mathrm{d}t}
    &&\text{（电场的环路定理：变化磁场激发电场）}\\
    &\oint_S\boldsymbol{B}\cdot\mathrm{d}\boldsymbol{S}=0
    &&\text{（磁场的高斯定理：磁场无源）}\\
    &\oint_L\boldsymbol{H}\cdot\mathrm{d}\boldsymbol{l}=\sum_{L\text{内}}I_{0i}+\frac{\mathrm{d}\Psi_D}{\mathrm{d}t}
    &&\text{（全电流定律：传导电流和变化电场激发磁场）}
    \end{align*}
    在使用时还需补充描述介质性质的三个关系式：
    \[
    \boldsymbol{D}=\varepsilon_0\varepsilon_r\boldsymbol{E},\qquad
    \boldsymbol{B}=\mu_0\mu_r\boldsymbol{H},\qquad
    \boldsymbol{j}=\sigma\boldsymbol{E}.
    \]

    \subsection{麦克斯韦方程组的物理意义}
    \begin{enumerate}
        \item 第一、三式分别是电场和磁场的高斯定理，说明电场有源（电荷是电场的源头）而磁场无源（不存在磁单极）；
        \item 第二式说明变化的磁场激发涡旋电场，这是电磁感应的普遍表述；
        \item 第四式说明传导电流和变化的电场都激发磁场，这是安培环路定理的推广；
        \item 四个方程把电与磁统一为一个整体——电磁场，预言了电磁波的存在，并指出光也是一种电磁波。麦克斯韦方程组是经典电磁学的最高概括，也是狭义相对论的重要实验与理论基础之一。
    \end{enumerate}

    \corepoint{四个方程中，第一式与第三式描述场的"源"，第二式与第四式描述场的"旋"。注意第二式的负号（楞次定律）和第四式中位移电流项的符号，二者共同保证电磁波能够自持传播。}

\section{平面电磁波}
    \subsection{电磁波的波动方程}
    在无源（$\rho=0$，$\boldsymbol{j}=0$）的均匀介质中，由麦克斯韦方程组可导出电场和磁场的波动方程
    \[
    \nabla^{2}\boldsymbol{E}=\frac{1}{u^{2}}\frac{\partial^{2}\boldsymbol{E}}{\partial t^{2}},\qquad
    \nabla^{2}\boldsymbol{B}=\frac{1}{u^{2}}\frac{\partial^{2}\boldsymbol{B}}{\partial t^{2}},
    \]
    其中
    \keyformula{u=\frac{1}{\sqrt{\varepsilon\mu}}=\frac{1}{\sqrt{\varepsilon_0\varepsilon_r\mu_0\mu_r}}=\frac{c}{\sqrt{\varepsilon_r\mu_r}}}
    为电磁波在介质中的传播速度。在真空中
    \[
    c=\frac{1}{\sqrt{\varepsilon_0\mu_0}}=2.998\times10^{8}\ \mathrm{m/s},
    \]
    与实测光速相符，这成为"光就是电磁波"的有力证据。

    平面电磁波沿 $x$ 轴传播时，其波函数可写为
    \[
    E=E_0\cos\omega\left(t-\frac{x}{u}\right),\qquad
    B=B_0\cos\omega\left(t-\frac{x}{u}\right).
    \]

    \subsection{平面电磁波的性质}
    由麦克斯韦方程组可导出平面电磁波的基本性质：
    \begin{enumerate}
        \item 电磁波是横波：$\boldsymbol{E}$、$\boldsymbol{B}$ 都与传播方向垂直，且 $\boldsymbol{E}\perp\boldsymbol{B}$；$\boldsymbol{E}\times\boldsymbol{B}$ 的方向即波的传播方向；
        \item $\boldsymbol{E}$ 与 $\boldsymbol{B}$ 同相位，即在空间任一点、任一时刻二者同时达到最大、同时为零；
        \item $\boldsymbol{E}$ 与 $\boldsymbol{B}$ 的量值满足 $E=uB$，在真空中 $E=cB$；
        \item 电磁波的传播速度由介质的性质决定，$u=c/\sqrt{\varepsilon_r\mu_r}$；
        \item 电磁波具有偏振性，是横波，故有偏振现象。
    \end{enumerate}

    \subsection{电磁波的能量和能流密度}
    电磁场中同时存在电场能量和磁场能量，其能量密度为
    \[
    w=w_e+w_m=\frac{1}{2}\varepsilon E^{2}+\frac{1}{2}\mu H^{2}.
    \]
    由于 $E=uB$、$\sqrt{\varepsilon}E=\sqrt{\mu}H$，可得 $w_e=w_m$，即电场能量与磁场能量各占一半，故
    \[
    w=\varepsilon E^{2}=\mu H^{2}.
    \]
    单位时间内通过垂直于传播方向的单位面积的能量，称能流密度（坡印廷矢量），为
    \keyformula{\boldsymbol{S}=\boldsymbol{E}\times\boldsymbol{H},\qquad S=EH=\frac{E^{2}}{\sqrt{\mu/\varepsilon}}}
    其平均值（波强）为
    \[
    \bar{S}=\frac{1}{2}E_0H_0=\frac{E_0^{2}}{2\sqrt{\mu_0/\varepsilon_0}}.
    \]

    电磁波谱按波长（或频率）排列，依次为：无线电波、微波、红外线、可见光、紫外线、$X$ 射线、$\gamma$ 射线，它们本质上都是电磁波，只是波长（频率）不同，因而表现出不同的产生方式和与物质相互作用的特性。

    \corepoint{（1）电磁波是横波，$\boldsymbol{E}\perp\boldsymbol{B}\perp\boldsymbol{u}$，三者构成右手螺旋；（2）$E=uB$ 是真空中 $E=cB$ 的推广，不能写成 $E=B$；（3）能流密度 $\boldsymbol{S}=\boldsymbol{E}\times\boldsymbol{H}$ 是瞬时值，实际测量中常用其时间平均值。}
